\chapter{Intégration et dérivation numérique}

Dans certains problèmes scientifiques, on arrive souvent à des dérivées
premières ou d'ordre $n$ ou à des intégrales simples ou doubles qui ne peuvent
pas être calculées par des méthodes analytiques. Il faut alors faire appel à
des méthodes numériques qui permettent de trouver des solutions approchées en
utilisant des formes discrètes issues des définitions mathématiques de base.

Ce chapitre présente les méthodes de dérivation numérique (en partant des
formules classiques, des interpolations linéaires, quadratiques et cubiques)
et des méthodes d'intégration numériques (formules des trapèzes et de
Simpson).

\section{Dérivation numérique}

\subsection{Formules classiques}

Soit $f$ une fonction différentiable sur un intervalle $[a,b]$. Le nombre
dérivé de $f$ au point $x$ est défini par :
\[
  f'(x) = \lim_{h\to 0}\frac{f(x+h) - f(x)}{h} .
\]
On peut également utiliser la formule
\[
  f'(x) = \lim_{h\to 0}\frac{f(x + h/2) - f(x - h/2)}{h} .
\]
De la même manière, la dérivée seconde est donnée par :
\[
  f''(x) = \lim_{h\to 0}\frac{f(x+h) - 2f(x) + f(x-h)}{h^2} .
\]
Numériquement, on écrit
\[
  f'(x) = \frac{f(x+h) - f(x)}{h}
  \qquad\text{ou}\qquad
  f'(x) = \frac{f\bigl(x + \frac{h}{2}\bigr) - f\bigl(x - \frac{h}{2}\bigr)}{h}
\]
et
\[
  f''(x) = \frac{f(x+h) - 2f(x) + f(x-h)}{h^2}
\]
où $h$ est un nombre très petit : par exemple $h = 10^{-3}$, $10^{-1}$.

La formule de la dérivée première peut s'obtenir à partir de la formule de
l'interpolation linéaire. En effet, la formule de l'interpolation linéaire est
\[
  f(x) = f(x_1)\frac{x - x_2}{x_1 - x_2} + f(x_2)\frac{x - x_1}{x_2 - x_1}
       + \frac{1}{2}f''(\xi)(x - x_1)(x - x_2) .
\]
D'où
\[
  f'(x) = \frac{f(x_2) - f(x_1)}{x_2 - x_1}
        + \frac{1}{2}f''(\xi)\bigl[(x - x_1) + (x - x_2)\bigr]
        + \frac{1}{2}(x - x_1)(x - x_2)\frac{\dd}{\dd x}\bigl[f''\bigl(\xi(x)\bigr)\bigr] .
\]
L'erreur de troncature correspondant à la dérivée par interpolation linéaire
est donc
\[
  E = \frac{1}{2}f''(\xi)\bigl[(x - x_1) + (x - x_2)\bigr]
    + \frac{1}{2}(x - x_1)(x - x_2)\frac{\dd}{\dd x}\bigl[f''\bigl(\xi(x)\bigr)\bigr]
\]
avec $\xi \in [x_1\,;\,x_2]$.

Si nous posons $x = x_1$ et $x_2 = x_1 + h$, on obtient
\[
  f'(x_1) = \frac{f(x_1 + h) - f(x_1)}{h}
\]
avec une erreur de troncature $E = -\dfrac{1}{2}hf''(\xi)$.

C'est la formule classique donnée précédemment.

\subsection{Formules plus précises}

\subsubsection{Dérivée à partir de l'interpolation quadratique}

La formule d'interpolation quadratique est
\[
  f(x) = f(x_1)f_1(x) + f(x_2)f_2(x) + f(x_3)f_3(x)
\]
où
\[
  f_1(x) = \frac{(x - x_2)(x - x_3)}{(x_1 - x_2)(x_1 - x_3)}, \quad
  f_2(x) = \frac{(x - x_1)(x - x_3)}{(x_2 - x_1)(x_2 - x_3)}, \quad
  f_3(x) = \frac{(x - x_1)(x - x_2)}{(x_3 - x_1)(x_3 - x_2)} .
\]
En procédant comme dans le cas d'une interpolation linéaire, et en posant
$x = x_2$, on établit que :
\[
  f'(x) = \frac{f(x+h) - f(x-h)}{2h}
\]
avec une erreur de troncature
\[
  E = -\frac{1}{6}h^2f'''(\xi) \qquad\text{où } \xi \in [x - h\,;\,x + h] .
\]
Si l'on posait $x = x_1$, il viendrait le résultat suivant :

Si $f$ est de classe $C^4$ sur l'intervalle $[a,b]$, alors pour toutes valeurs
de $x$ et $h$ avec $a < x + 2h < b$, il existe $\xi \in [x\,;\,x + 2h]$ tel
que
\[
  f'(x) = \frac{1}{2h}\bigl[-3f(x) + 4f(x+h) - f(x+2h)\bigr] + \frac{1}{3}h^2f'''(\xi) .
\]

\subsubsection{Dérivée à partir de l'interpolation cubique}

On obtient le résultat suivant :

Si $f$ est de classe $C^6$ sur l'intervalle $[a,b]$ alors pour toute valeur de
$x$ et $h$ telles que $a < x - 2h < x + 2h < b$, il existe
$\xi \in [x - 2h\,;\,x + 2h]$ tel que
\[
  f'(x) = \frac{1}{12h}\bigl[f(x-2h) - 8f(x-h) + 8f(x+h) - f(x+2h)\bigr]
        + \frac{1}{30}h^4f^{(5)}(\xi) .
\]

\begin{exercice}[Exercice 1]
Établir la formule de dérivation numérique à partir de l'interpolation
quadratique.
\end{exercice}

\begin{exercice}[Exercice 2]
Au cours d'une expérience de dynamique, la position en fonction du temps est
donnée dans le tableau suivant :
\begin{center}
  \begin{tabular}{|c|c|c|c|c|c|c|}
    \hline
    $t$ & 1.0 & 1.10 & 1.20 & 1.30 & 1.40 & 1.50\\ \hline
    $x$ & 1.6487 & 1.7333 & 1.8221 & 1.9155 & 2.0138 & 2.1170\\ \hline
  \end{tabular}
\end{center}
Calculer les vitesses aux instants $t = 1.05$ ; $1.15$ ; $1.25$ ; $1.35$ et
$1.45$.
\end{exercice}

\section{Intégrations numériques -- Formules de Newton-Cotes}

\subsection*{L'intégrale définie}

\[
  I = \int_a^b f(x)\,\dd x
\]
représente un nombre qui, dans certains cas favorables peut être calculé
analytiquement. Dans le cas contraire, on doit faire appel au calcul
numérique. Pour ce faire, la fonction $f$ supposée continue sur $[a,b]$ est
remplacée par une formule d'interpolation. Généralement, $f$ est approximée
par un polynôme et les formules d'intégration qui s'en déduisent sont dites
\emph{formules de Newton-Cotes}. Ces formules supposent que les n\oe uds
d'interpolation sont équidistants.

\subsection{Méthode des trapèzes}

\subsubsection{Formule composite des trapèzes}

Dans la méthode des trapèzes, l'interpolation de la fonction $f$ est
linéaire. En supposant que l'intervalle $[a,b]$ est court, on a
\[
  f(x) = f(a)\frac{x - b}{a - b} + f(b)\frac{x - a}{b - a}
       + \frac{1}{2}f''(\xi)(x - a)(x - b) .
\]
Il vient alors
\[
  I \cong \frac{b - a}{2}\bigl[f(a) + f(b)\bigr] + E
\]
où $E$ est l'erreur de troncature définie par
\[
  E = \frac{1}{2}\int_a^b f''(\xi)(x - a)(x - b)\,\dd x .
\]
Pour évaluer $E$, nous utilisons le théorème de la moyenne pour les
intégrales.

\paragraph{Théorème de la moyenne} Soient $f$ et $g$ deux fonctions
continues sur $[a,b]$. Si $g(x)$ ne change pas de signe sur $[a,b]$, alors il
existe $\eta \in [a,b]$ tel que
\[
  \int_a^b f(x)g(x)\,\dd x = f(\eta)\int_a^b g(x)\,\dd x .
\]
Puisque le produit $(x - a)(x - b)$ reste négatif sur $[a,b]$, nous aurons
donc
\[
  E = -\frac{1}{12}f''(\xi)(b - a)^3 .
\]
Si $b - a = h$, alors
\[
  E = -\frac{1}{12}f''(\xi)h^3 .
\]
Lorsque la longueur de l'intervalle $[a,b]$ n'est pas assez petite, l'on
utilise la propriété d'additivité de l'intégrale pour faire un découpage en
morceaux réduits, à savoir :
\[
  \int_a^b f(x)\,\dd x = \sum_{k=1}^{n}\int_{x_{k-1}}^{x_k} f(x)\,\dd x
  \qquad\text{où } a = x_0 < x_1 < \dots < x_{n-1} < x_n = b .
\]
Prenons $x_k = x_0 + kh$, on obtient alors :
\[
  \int_a^b f(x)\,\dd x = \frac{h}{2}\bigl[f(x_0) + 2f(x_1) + 2f(x_2) + \dots
                         + 2f(x_{n-1}) + f(x_n)\bigr] + E
\]
où $E = \sum_{k=1}^{n} E_k$.

C'est la formule composite de la méthode des trapèzes.

Sachant que $E_k = -\dfrac{1}{12}h^3f''(\xi_k)$, on obtient
\[
  E = \sum_{k=1}^{n} E_k = -\frac{1}{12}h^3\sum_{k=1}^{n} f''(\xi_k)
    = -\frac{nh^3}{12}\,\frac{\sum_{k=1}^{n} f''(\xi_k)}{n}
    = -\frac{nh^3}{12}f''(\eta)
\]
où $\eta \in [a,b]$. Or $b - a = nh$. D'où
\[
  E = -\frac{h^2}{12}(b - a)f''(\eta) .
\]

\begin{exercice}[Exercice 3]
Calculer à la main $I = \displaystyle\int_1^2 \frac{\dd x}{x}$ pour $h = 0.2$,
puis pour $h = 0.1$. Comparer à la valeur exacte.
\end{exercice}

\subsubsection{Algorithme de la formule composite des trapèzes}

Cet algorithme permet de calculer une intégrale par la méthode composite des
trapèzes. Il utilise des tableaux pour $x_k$ et $f_k$.

\begin{algo}
  \State 1\textdegree- Définir la fonction $f$
  \State 2\textdegree- Donner $a$, $b$ et $n$
  \State 3\textdegree- $h \leftarrow \dfrac{b - a}{n}$, $S \leftarrow 0$
  \State 4\textdegree- \textbf{Pour} $k$ allant de $0$ jusqu'à $n$, \textbf{faire}
  \State \hspace{2.5em} $x_k \leftarrow a + kh$
  \State \hspace{2.5em} $f_k \leftarrow f(x_k)$
  \State \hspace{1.2em} \textbf{Fin pour}
  \State 5\textdegree- \textbf{Pour} $k$ allant de $1$ jusqu'à $n-1$, \textbf{faire}
  \State \hspace{2.5em} $S \leftarrow S + hf_k$
  \State \hspace{1.2em} \textbf{Fin pour}
  \State \hspace{1.2em} $S \leftarrow S + \dfrac{h}{2}(f_0 + f_n)$
  \State 6\textdegree- \textbf{Écrire} $S$
\end{algo}

\begin{exercice}[Exercice 4]
Écrire un autre algorithme qui n'utilise pas les tableaux.
\end{exercice}

\subsection{Méthode de Simpson}

Dans la méthode de Simpson, la fonction $f$ est approximée par un polynôme de
degré deux qui passe par les points $a$, $b$ et $c = \dfrac{a+b}{2}$. Cette
méthode permet d'obtenir :
\[
  \int_a^b f(x)\,\dd x = \frac{b - a}{6}\bigl[f(a) + 4f(c) + f(b)\bigr] + E
\]
avec
\[
  E = \frac{1}{6}\int_a^b f'''(\xi)(x - a)(x - b)(x - c)\,\dd x .
\]
En utilisant l'additivité des intégrales, on obtient :
\[
  \int_a^b f(x)\,\dd x = \frac{h}{3}\bigl[f(x_0) + 4f(x_1) + 2f(x_2) + 4f(x_3)
                         + 2f(x_4) + \dots + 4f(x_{n-1}) + f(x_n)\bigr] .
\]

\begin{exercice}[Exercice 5]
Écrire un algorithme pour la méthode de Simpson.
\end{exercice}

\begin{exercice}[Exercice 6]
Établir une formule d'intégration numérique qui utilise une interpolation de
degré 3 et qui interpole les points $a$, $b$, $c = \dfrac{a+b}{4}$ et
$d = \dfrac{3(a+b)}{4}$.
\end{exercice}

\begin{exercice}[Exercice 7]
Utiliser la méthode de Simpson pour calculer manuellement l'intégrale
$I = \displaystyle\int_0^1 \frac{4}{1+x^2}\,\dd x$ avec $h = 0.1$.
\end{exercice}

\section{Intégrales mixtes et doubles}

\subsection{Intégrale mixte}

C'est une intégrale de la forme
\[
  I = \int_a^b f(x)g(x)\,\dd x
\]
où l'expression mathématique de $g(x)$ est connue et où $f(x)$ est approchée
par une formule d'interpolation. Ce type d'intégrale intervient par exemple
dans le calcul des coefficients de Fourier du développement d'une fonction :
\[
  a_n = \frac{2}{T}\int_0^T f(x)\cos(nx)\,\dd x
  \qquad\text{et}\qquad
  b_n = \frac{2}{T}\int_0^T f(x)\sin(nx)\,\dd x .
\]
Dans ce cas $g(x) = \cos(nx)$ ou $\sin(nx)$. Ce type d'intégrale intervient
également dans la résolution des équations intégrales (par exemple
l'équation $f(x) = 5x + \int_0^1 (t + x)f(t)\,\dd t$).

Pour calculer ce type d'intégrale, on élabore des formules d'intégration de
type Newton-Cotes ($f(x)$ est interpolée par un polynôme).

\subsection{Interpolation double}

Il s'agit de calculer une intégrale de la forme
\[
  I = \int_{x=a}^{x=b}\int_{y=c}^{y=d} Z(x,y)\,\dd x\,\dd y .
\]
Pour ce faire, l'on peut utiliser la formule d'interpolation double de la
fonction $Z(x,y)$.
